\subsection*{一步一步建立直观}

\paragraph{1. 先把累积当作一台计数器}

把横轴看作时间，把曲线高度 $x^2$ 看作某个时刻的流入速率。一个窄矩形的宽是持续时间，高是近似速率，因此面积是这段时间流入的量。先只放少量矩形：每块都把一段变化中的速率当成常数，所以结果会偏离真正的累积。



\paragraph{2. 切得更细，究竟改善了什么}

增加分割数时，每段内速率的变化变小，矩形顶部更贴近曲线。这里函数连续且递增，左端点和低估面积。关键不是图形越来越好看，而是可以把每段误差相加并控制总误差；无穷多小块只能在这样一个极限过程中得到明确含义。



\paragraph{3. 只移动终点，寻找新增的量}

保持分割足够细，再改变右端点 $b$。原先已经累计的部分没有改变，新增部分是一条宽度约为 $h$、高度约为 $b^2$ 的窄带。因此面积函数的增量约为 $b^2h$。将增量除以 $h$ 后取极限，就把“累积的变化率”接回了原速率。


